\documentclass{ctexart}
\usepackage{enumitem}
\usepackage{graphicx}
\usepackage{float}
\usepackage{amsmath,amssymb}
\usepackage[ruled,linesnumbered]{algorithm2e}
\usepackage[a4paper,margin=2cm]{geometry}
\title{\vspace*{-1cm}Algorithm Design and Analysis Homework 5}
\author{Tianyi Guan 2400017423}
\date{\today}

\begin{document}
\maketitle

\begin{enumerate}[label=\arabic*.]
    \item \textbf{Problem 7.4}\\
        由于 $N$ 是简单容量网络，由定义，其各边容量均为 $1$。而 $f$ 也是 0-1 可行流，意味着每条边上的流量必为 0 或 1。因此，考虑 $N(f)$ 每条边的计算过程，所有的 $N(f)$ 上的正向边，其容量只可能是 0（当边容量为 1 且真实流量为 1，此时正向边消失）或 1（当边容量为 1 且真实流量为 0）；对于所有 $N(f)$ 上的反向边，其容量只可能是 0（真实流量为 0，此时反向边不生成）或 1（真实流量为 1）。这就证明了 $N(f)$ 中所有存在的边，容量必然均为 1。
        接下来证明节点结构的限制。对于 $N$ 中的任意中间节点 $v$，其原本满足入度为 1 或出度为 1。
        我们以 $v$ 的入度为 1，出度为 $k$ 为例：
        \begin{itemize}
            \item \textbf{当经过 $v$ 的流量为 0 时：} $N(f)$ 中与 $v$ 关联的边方向不发生任何改变，因此在 $N(f)$ 中 $v$ 的入度仍为 1。
            \item \textbf{当经过 $v$ 的流量为 1 时：} 上游唯一的那条入边必定满流，在 $N(f)$ 中方向反转，转变为 $v$ 的一条出边；而在下游的 $k$ 条出边中，必然有且只有一条边满流，在 $N(f)$ 中方向反转，转变为 $v$ 的一条入边。其余 $k-1$ 条出边流量为 0，在 $N(f)$ 中依然保持为出边。此时清点 $N(f)$ 中 $v$ 的度数可知，其入度仅有那条由反向边生成的入边，总入度依然严格保持为 1。
        \end{itemize}
        同理可证，若原网络中 $v$ 的入度为 $k$，出度为 1，无论是否有流量经过，经过类似的边反向推导可知，其在 $N(f)$ 中的出度永远保持为 1。

        综上所述，在辅助网络 $N(f)$ 中，各边容量均为 1，且所有的中间点依然严格满足或者入度为 1 或者出度为 1 的结构性质。因此，$N(f)$ 仍然是简单容量网络。
    \item \textbf{Problem 7.6}\\
        只需要在原有的$N$上加入两个超节点$S*$和$T*$，其中$S*$通过所有出边指向原本的所有发点，$T*$通过所有入边连接原本的所有收点，且所有边容量都为$+\infty$。剩下对这个扩充网络解最大流问题就可以了。
    \item \textbf{Problem 7.7}\\
        只需要使用“拆点法”。
        对于原网络 $N$ 中的每一个带有顶点容量限制的中间节点 $u \in V - \{s, t\}$，执行以下操作：
        将节点 $u$ 拆分成两个全新的节点：入节点 $u_{in}$ 和出节点 $u_{out}$。
        在 $u_{in}$ 和 $u_{out}$ 之间新增一条有向边 $\langle u_{in}, u_{out} \rangle$，并将其容量设为原顶点的容量，即 $c(u_{in}, u_{out}) = c(u)$。
        将原网络中所有指向 $u$ 的入边全部改接到 $u_{in}$ 上；所有从 $u$ 出发的出边全部改由 $u_{out}$ 引出。所有原边的容量均保持不变。
        转化完成后，任何流经原节点 $u$ 的流量都必须且只能经过新增的内部边 $\langle u_{in}, u_{out} \rangle$，从而完美受到 $c(u)$ 的物理限制。剩下只需要对这个扩充后的新网络解标准的最大流问题就可以了。
    \item \textbf{Problem 7.17}
        \begin{enumerate}[label=(\arabic*)]
            \item 一个简单的想法是直接先调用dinic算法求出最大流，然后以这个求出的最大流量作为输入，用floyd算法在辅助网络上找负权环路，使得费用最小化。
                但一直在辅助网络上消圈效率比较低，因为每次floyd算法都会带来$O(V^3)$的时间开销。
                考虑如下的优化算法：\\
                \begin{algorithm}[H]
                \caption{连续最短路算法 (Successive Shortest Path, SSP)}
                \label{algo:ssp}
                \KwIn{带费用的容量网络 $N = (V, E, c, w, s, t)$，其中 $c$ 为容量，$w$ 为单位费用}
                \KwOut{从源点 $s$ 到汇点 $t$ 的最小费用最大流 $F$ 及最小总费用 $C$}
                
                \tcc{1. 初始化}
                $F \leftarrow 0$\;
                $C \leftarrow 0$\;
                $f(u,v) \leftarrow 0, \quad \forall \langle u,v \rangle \in E$\;
                
                \tcc{2. 核心贪心推流主循环}
                \While{\text{True}}{
                    根据当前流量 $f$，构造辅助网络 $N(f)$\;
                    \tcc{对于 $N(f)$ 中的边：正向边费用为 $w$，反向边费用为 $-w$}
                    
                    使用 SPFA/Bellman-Ford 在 $N(f)$ 中寻找从 $s$ 到 $t$ 的最小费用最短路径 $P$\;
                    
                    \eIf{找不到路径 $P$}{
                        \textbf{break}\; \tcp*{无法继续推流，已达到最大流}
                    }{
                        \tcc{3. 计算瓶颈与推流}
                        $\Delta f \leftarrow \min_{\langle u,v \rangle \in P} \{ c_f(u,v) \}$\; \tcp*{找出路径 $P$ 上的最小残余容量}
                        
                        $F \leftarrow F + \Delta f$\;
                        $C \leftarrow C + \Delta f \times \sum_{\langle u,v \rangle \in P} w(u,v)$\;
                        
                        \tcc{4. 更新实际流量}
                        \ForEach{边 $\langle u,v \rangle \in P$}{
                            \eIf{$\langle u,v \rangle$ 是原图中的正向边}{
                                $f(u,v) \leftarrow f(u,v) + \Delta f$\;
                            }{
                                $f(v,u) \leftarrow f(v,u) - \Delta f$\; \tcp*{利用反向边退流}
                            }
                        }
                    }
                }
                
                \Return 最大流量 $F$, 最小费用 $C$\;
            \end{algorithm}
        \item 只需要稍微修改上面的ssp算法，算法伪码如下：\\
            \begin{algorithm}[H]
            \caption{带预算限制的连续最短路算法 (Budget-Constrained SSP)}
            \label{algo:budget_ssp}
            \KwIn{容量网络 $N$，边容量 $c$，边费用 $w$，源点 $s$，汇点 $t$，\textbf{总预算 $B$}}
            \KwOut{在费用不超过 $B$ 的前提下的最大流 $F$}
            
            $F \leftarrow 0$\;
            $C \leftarrow 0$\;
            $f(u,v) \leftarrow 0, \quad \forall \langle u,v \rangle \in E$\;
            
            \While{\text{True}}{
                根据当前流量 $f$，构造辅助网络 $N(f)$\;
                使用 SPFA 寻找 $N(f)$ 中从 $s$ 到 $t$ 的最小费用最短路径 $P$，设其单位流量费用为 $w_P$\;
                
                \eIf{找不到路径 $P$}{
                    \textbf{break}\; \tcp*{网络已达物理极限，提前结束}
                }{
                    $\Delta f \leftarrow \min_{\langle u,v \rangle \in P} \{ c_f(u,v) \}$\; \tcp*{计算物理瓶颈}
                    $B_{\text{remain}} \leftarrow B - C$\; \tcp*{计算剩余预算}
                    
                    \eIf{$\Delta f \times w_P \le B_{\text{remain}}$}{
                        \tcc{预算充足，拉满物理瓶颈}
                        $f_{\text{push}} \leftarrow \Delta f$\;
                    }{
                        \tcc{预算不足，提前刹车（假设允许非整数流，若不可拆分则向下取整）}
                        $f_{\text{push}} \leftarrow B_{\text{remain}} / w_P$\;
                    }
                    
                    \tcc{执行推流、记账与残余网络更新}
                    $F \leftarrow F + f_{\text{push}}$\;
                    $C \leftarrow C + f_{\text{push}} \times w_P$\;
                    更新路径 $P$ 上所有正反向边的流量 $f$\;
                    
                    \If{$f_{\text{push}} < \Delta f$}{
                        \textbf{break}\; \tcp*{预算已耗尽，终止推流}
                    }
                }
            }  
            \Return 最大可用流量 $F$\;
        \end{algorithm}
        \end{enumerate}
    \item \textbf{Problem 7.21 (1)}\\
        \textbf{1. 节点定义：}
        \begin{itemize}
            \item 设立一个超级源点 $S$ 和一个超级汇点 $T$。
            \item 设立四个中间节点 $Q_1, Q_2, Q_3, Q_4$，分别代表四个季度的仓库。
        \end{itemize}

        \textbf{2. 边集与参数定义（容量 $c$, 单位费用 $w$）：}

        \textbf{第一类：生产边（从源点到各季度，代表当季生产能力与成本）}
        从 $S$ 向每个 $Q_i$ 连边，容量为当季最大生产能力，费用为生产成本：
        \begin{itemize}
            \item $\langle S, Q_1 \rangle$: $c = 25, w = 12.0$
            \item $\langle S, Q_2 \rangle$: $c = 35, w = 11.0$
            \item $\langle S, Q_3 \rangle$: $c = 30, w = 11.5$
            \item $\langle S, Q_4 \rangle$: $c = 20, w = 12.5$
        \end{itemize}

        \textbf{第二类：库存边（相邻季度之间，代表跨季仓储）}
        从 $Q_i$ 向 $Q_{i+1}$ 连边，代表本季未交货的产品存入仓库留至下季，容量无限制，费用为库存费：
        \begin{itemize}
            \item $\langle Q_1, Q_2 \rangle$: $c = +\infty, w = 0.1$
            \item $\langle Q_2, Q_3 \rangle$: $c = +\infty, w = 0.1$
            \item $\langle Q_3, Q_4 \rangle$: $c = +\infty, w = 0.1$
        \end{itemize}

        \textbf{第三类：交货边（从各季度到汇点，代表满足当季合同需求）}
        从每个 $Q_i$ 向 $T$ 连边，容量为当季必须交货的数量，由于交货本身不产生额外内部费用，费用设为 0：
        \begin{itemize}
            \item $\langle Q_1, T \rangle$: $c = 15, w = 0$
            \item $\langle Q_2, T \rangle$: $c = 20, w = 0$
            \item $\langle Q_3, T \rangle$: $c = 25, w = 0$
            \item $\langle Q_4, T \rangle$: $c = 20, w = 0$
        \end{itemize}

        \textbf{3. 模型求解：}
        只需要在这个构建好的网络中，求解从 $S$ 到 $T$ 的\textbf{总流量为 80 的最小费用流}（或者直接求最小费用最大流，由于汇点入边总容量恰好为 80，最大流必为 80）。求得的最小费用即为最低总成本，每条边上的实际流量即为最优的生产与库存调度计划。
\end{enumerate}

\end{document}